\documentclass[12pt]{article}

% Packages for math, formatting, and layout
\usepackage{amsmath}
\usepackage{amssymb}
\usepackage{geometry}
\usepackage{hyperref}
\usepackage{enumitem}

% Set page margins
\geometry{a4paper, margin=1in}

\title{\textbf{Value Function Approximation and DQN}}
\author{}
\date{}

\begin{document}

\maketitle

\subsection*{Why Tabular Methods Fail?}
\begin{itemize}[label=$\bullet$]
    \item Up until now, we have studied algorithms like standard Q-Learning and SARSA, which rely on a Q-table. A Q-table requires a distinct row for every possible state.
    \item The problem with these tabular methods is that the number of possible states can be astronomically large, making it impossible to store a complete table.
    \item To solve this problem, we must move from \textit{memorization} to \textit{generalization}.
    \item The agent should not repeatedly learn from scratch for states with similar surroundings; it should generalize its knowledge across similar states.
\end{itemize}

\subsection*{Value Function Approximation}
\begin{itemize}[label=$\bullet$]
    \item Instead of using a lookup table, we represent the Q-value as a parameterized mathematical function:
    $$Q(s, a) \approx Q(s, a; \theta)$$
    where $\theta$ represents the parameters (or weights) of our function.
\end{itemize}

\subsection*{Linear vs. Non-Linear Approximators}
\begin{enumerate}
    \item \textbf{Linear Approximators:} We manually extract features from the state ($\phi(s)$) and multiply them by our weights: 
    $$Q(s, a; \theta) = \theta^\top \phi(s, a)$$
    \begin{itemize}[label=$\bullet$]
        \item This approach is mathematically stable and guaranteed to converge.
    \end{itemize}
    \item \textbf{Non-Linear Approximators (Neural Networks):} We pass the raw state into a Deep Neural Network (DNN), and it outputs the Q-values.
    \begin{itemize}[label=$\bullet$]
        \item Here, the network learns and extracts the features on its own.
    \end{itemize}
\end{enumerate}

\subsection*{Gradient Descent in RL}
\begin{itemize}[label=$\bullet$]
    \item To train our approximator, we need to update our weights $\theta$ to minimize our errors. We use Gradient Descent for this. 
    \item We define a loss function (which measures how far off our prediction is from the target) and update $\theta$ in the direction that reduces this loss:
    $$\theta \leftarrow \theta - \alpha \nabla_\theta \text{Loss}(\theta)$$
\end{itemize}

\subsection*{The DQN Loss Function}
\begin{itemize}[label=$\bullet$]
    \item Deep Q-Networks (DQN) train the neural network by minimizing the Mean Squared Error (MSE) between the network's current Q-value prediction and the Bellman TD Target.
    \item The loss function $L(\theta)$ at iteration $i$ is defined as:
    $$L_i(\theta_i) = \mathbb{E}_{(s, a, r, s') \sim U(D)} \left[ \left( r + \gamma \max_{a'} Q(s', a'; \theta_i^-) - Q(s, a; \theta_i) \right)^2 \right]$$
    \item \textbf{Where:}
    \begin{itemize}[label=$\circ$]
        \item $\theta_i$: The weights of the primary Q-network we are currently training.
        \item $\theta_i^-$: The weights of the separate Target Network.
        \item $U(D)$: Random uniform sampling from the Experience Replay Buffer ($D$).
        \item $r + \gamma \max_{a'} Q(s', a'; \theta_i^-)$: The TD Target (what we want our network to output).
        \item $Q(s, a; \theta_i)$: The current prediction from our primary network.
    \end{itemize}
\end{itemize}

\subsection*{Why Experience Replay Helps?}
\begin{itemize}[label=$\bullet$]
    \item In RL, an agent generally learns from consecutive frames of an episode.
    \item This data is highly correlated.
    \item If you feed highly correlated, sequential data into a neural network, it will overfit to that specific sequence and fail to generalize.
    \item \textbf{Solution:} The agent stores all its experiences $(s, a, r, s')$ in a massive dataset called the \textit{Experience Replay Buffer}. 
    \item During training, the network pulls random mini-batches of data from this buffer. This breaks the temporal correlation of the data, satisfying the neural network's requirement for independent and identically distributed (i.i.d.) data.
\end{itemize}

\subsection*{Why Target Networks Help?}
\begin{itemize}[label=$\bullet$]
    \item In standard Q-learning, the agent updates its value estimate based on its estimate of the next state.
    \item If you use the exact same neural network to calculate both the prediction and the target, the target shifts every single time the network's weights update.
    \item This is known as the \textbf{"moving target" problem}, and it causes feedback loops that make the network spiral out of control.
    \item \textbf{Solution:} DQN splits the architecture into two separate networks:
    \begin{enumerate}
        \item The \textbf{primary network} ($\theta$) is updated continuously.
        \item The \textbf{Target Network} ($\theta^-$) is a frozen clone of the primary network, used only to calculate the TD target.
    \end{enumerate}
    \item Because of this separation, the targets remain stationary and stable while the primary network learns.
\end{itemize}

\end{document}